\documentclass[openany,fontset=fandol]{ctexbook}
\usepackage{amsmath,amssymb}
\usepackage[unicode,hidelinks]{hyperref}
\newcommand{\note}[1]{\footnote{译注：#1}}
\begin{document}
\title{代数几何基础第五卷预备笔记\\第12节：线性截面族与爆破说明}
\author{据用户上传英译编辑稿译出}
\date{2026年11月17日}
\maketitle
\chapter*{范围与校订状态}
本批连续译出底本第48页前五段，自“Once we grant this”起，止于“the subscheme $Z^1$ corresponding to $X^{(m)}$ supersingular is nowhere dense”。其中末一长段按论述层次分段排印。下一接点是同页“It remains, therefore, to prove a lemma”。

正文保留原作者的推广计划，不补成独立定理证明。底本在基域、纤维下标及所用记号上存在残缺或混用，均在相应位置注明；不能把这些问题视为已由翻译消除。此前缺文继续保留。本稿独立保存，未覆盖或合入主稿。
\chapter*{§12 将上述结果推广到线性截面（续）}
一旦承认上述结果，就得到推广8.8、8.9、8.11的那些显然推论。无疑，也可以在任意$m$的情形下陈述第8节的其他命题。\note{“承认上述结果”承接前批仍待验证的推广讨论，不表示本批已经补齐证明；尤其不能据此恢复8.8等处的既有缺文。}

如果这需要额外的编写工作，那么即使收入前面的微分方面的讨论，我们也可以放弃这一推广。

第9节的结果也是如此。

至于第10节，其中研究的情形按如下方式推广到任意$m$。固定$P$中一个余维为$m+1$的线性子簇$C$，考虑由$P$中包含$C$的余维$m$线性子簇$L$构成的射影空间$Q$。$Q$是$\operatorname{Grass}^m$的闭子概形；特别地，可以构造我们打算研究的$X_Q^{(m)}$。\note{这里沿用本节$X^{(m)}$在格拉斯曼参数空间上的普遍线性截面族，$X_Q^{(m)}$表示将该族限制到$Q$。底本的“close subscheme”按语境译作“闭子概形”。}

\newpage
第一点无论如何都应当写入正文：至少当$C$与$X$“正则地”相交时，态射
\[
X_Q^{(m)}\longrightarrow X
\]
仍是双有理的。更确切地说，在这种情况下，$X_Q^{(m)}$典范地同构于沿$X\times_P C$爆破$X$所得到的预概形。这一事实的证明，无非是借助10.3运用10.2。\note{底本给“regularly”加引号而未在本段重述精确定义；本译文保留这一限定，不把爆破同构或双有理性扩展为任意$C$均成立的结论。}

第二点颇有意义，但并非一定要收入：若把$C$选得“足够一般”，则
\[
X_Q^{(m)}\longrightarrow Q
\]
具有某些良好的性质。其中最经典的一条如下。假设$X$在$S=\operatorname{Spec}$〔底本此处未写出所取谱的对象〕上光滑，维数为$n\geq m$，并且固有、几何不可约。对于“足够一般”的$C$，令$T$为那些$\xi\in Q$组成的集合，使得$X_\xi^{(m)}$不是$k(\xi)$上维数为$n-m$的光滑概形。则$T$几何不可约〔底本此处另写“在$K(\xi)$上”，所指基域待核〕，且在$Q$中余维为一。令$T'$为$T$中那些$\xi$组成的集合，使得$X_\xi^{(m)}$至少有一个“超奇异”点，则$T'$在$T$中无处稠密。\note{底本前两处只写$X^{(m)}$，但同段后文明确写$X_\xi^{(m)}$，并逐点使用剩余域，故正文补明纤维下标。底本大小写混用$K(\xi)$、$k(\xi)$，纤维处统一作$k(\xi)$。但全局判别集合$T$的几何不可约性究竟相对于哪个基域，不能从“$S=\operatorname{Spec}$”的残句和游移的$K(\xi)$唯一恢复，故不擅自补定。}

最后，若$f\colon X\to P$是浸入，则略微扩大$T'$之后，对每个$\xi\in T\setminus T'$，$X_\xi^{(m)}$中恰有一个非光滑点，且该点在$k(\xi)$上有理。我忘了在陈述中说明：假设$X\to P$未分歧，而且一开始必须把$f$替换为$\phi_n f$，其中$n\geq2$（$\phi_n$在9.1中定义）。\note{底本在“浸入”处写大写$F$，而本段随后仍以$f$表示同一态射，正文统一为$f$。底本又用$n$既表示$X$的维数，又作为$\phi_n$的指标；这里照录这一重用并指出，不自行断言两个用途必须取同一数值。作者后补的未分歧假设及预先替换$f$的要求仍是陈述的一部分，不能遗漏。}

证明这一陈述最自然的办法，似乎是利用$\operatorname{Grass}^m$中使线性截面“奇异”的子概形$Z$（在8.8中记为$T$）。我们看到，在所给条件下，它几何不可约，余维为一，而对应于线性截面超奇异的子概形$Z^1$无处稠密。\note{底本此句交替使用$X^m$和$X^{(m)}$而未写参数下标；正文以“线性截面”表达其参数条件，原记号保留于此。本句仍是原作者提出的证明路径，不代表本批已经核实判别子概形的全部性质。}
\end{document}
